\documentclass[a4paper,12pt]{article}
\usepackage[utf8]{inputenc}
\usepackage{geometry}
\usepackage{amsmath, amssymb}
\usepackage{booktabs}
\usepackage{hyperref}
\usepackage[T5]{fontenc}
\usepackage[vietnamese, english]{babel}

\geometry{margin=1in}

\title{Searching and Sorting Algorithms Report}
\author{Nguyễn Dương Phúc}
\date{10/5/2025}

\begin{document}
\maketitle
\tableofcontents
\newpage

\section{Introduction}
This report explains and analyzes common \textbf{searching} and \textbf{sorting} algorithms using simple examples.  

We use the unsorted array:
\[
[34, 7, 23, 32, 5, 62, 36, 14]
\]
for sorting, and the sorted array:
\[
[5, 7, 14, 23, 32, 34, 36, 62]
\]
for searching.

\newpage
\section{Searching Algorithms}

\subsection{Linear Search}
\textbf{Idea:} Check each element one by one until the target is found.  \\
\textbf{Steps (find 36):}
\begin{itemize}
  \item Step 1: Compare 5 → not target
  \item Step 2: Compare 7 → not target
  \item Step 3: Compare 14 → not target
  \item Step 4: Compare 23 → not target
  \item Step 5: Compare 32 → not target
  \item Step 6: Compare 34 → not target
  \item Step 7: Compare 36 → target found at index 6
\end{itemize}
\textbf{Time complexity:} O(n), space = O(1), best-case = O(1), worst-case = O(n).  \\
\textbf{Advantages:} Works on sorted or unsorted arrays; simple; no extra memory needed. \\
\textbf{Disadvantages:} Slow for large arrays.  \\
\textbf{When to use:} Small datasets, contiguous memory.

\subsection{Binary Search}
\textbf{Idea:} Works on sorted arrays. Pick middle element and reduce search space by half.  
\textbf{Steps (find 36):}
\begin{itemize}
  \item Step 1: left=0, right=7 → mid=(0+7)/2=3 → arr[3]=23 → 36>23 → left=4
  \item Step 2: left=4, right=7 → mid=(4+7)/2=5 → arr[5]=34 → 36>34 → left=6
  \item Step 3: left=6, right=7 → mid=(6+7)/2=6 → arr[6]=36 → target found
\end{itemize}
\textbf{Time complexity:} O(log n), space=O(1), best-case=O(1), worst-case=O(log n).  \\
\textbf{Advantages:} Very fast for large sorted arrays; minimal memory.\\
\textbf{Disadvantages:} Requires sorted array; not suitable for unsorted data.\\
\textbf{When to use:} Large, sorted arrays.

\subsection{Jump Search}
\textbf{Idea:} Jump in blocks of size $\sqrt{n}$, then do linear search in that block.  \\
\textbf{Steps:}
\begin{itemize}
  \item Step 1: block size = 2 → jump arr[1]=7 → 36>7
  \item Step 2: jump arr[3]=23 → 36>23
  \item Step 3: jump arr[5]=34 → 36>34
  \item Step 4: jump arr[7]=62 → 36<62 → search block arr[6:7]
  \item Step 5: arr[6]=36 → target found
\end{itemize}
\textbf{Time complexity:} O($\sqrt{n}$), space=O(1).\\
\textbf{Advantages:} Faster than linear search for large sorted arrays; simple to implement.\\
\textbf{Disadvantages:} Requires sorted array; block size affects performance.\\
\textbf{When to use:} Large sorted arrays where binary search is not practical.

\subsection{Ternary Search}
\textbf{Idea:} Split array into 3 parts, compare with 2 midpoints.  \\
\textbf{Steps:}
\begin{itemize}
  \item Step 1: l=0, r=7 → mid1=2, mid2=5 → arr[2]=14, arr[5]=34 → 36>34 → search right third
  \item Step 2: l=6, r=7 → mid1=6, mid2=7 → arr[6]=36 → target found
\end{itemize}
\textbf{Time complexity:} O(log$_3$ n), space=O(1)\\
\textbf{Advantages:} Similar to binary search; theoretically reduces comparisons.\\
\textbf{Disadvantages:} Only works on sorted arrays; slightly more complex than binary search.\\
\textbf{When to use:} Sorted arrays; sometimes educational purposes.

\subsection{Interpolation Search}
\textbf{Idea:} Estimate position using formula: $pos = low + \frac{(x-arr[low])*(high-low)}{arr[high]-arr[low]}$  \\
\textbf{Steps:} pos = 6 → arr[6]=36 → found  \\
\textbf{Time complexity:} O(log log n) for uniform data, worst-case O(n)\\
\textbf{Advantages:} Very fast for uniformly distributed sorted data.\\
\textbf{Disadvantages:} Poor performance for non-uniform data; requires sorted array.\\
\textbf{When to use:} Sorted, uniformly distributed data.
\subsection{Exponential Search}
\textbf{Idea:} Check elements exponentially until overshoot, then binary search.  \\
\textbf{Steps:}
\begin{itemize}
  \item Step 1: i=0 → arr[0]=5 → 36>5
  \item Step 2: i=1 → arr[1]=7 → 36>7
  \item Step 3: i=2 → arr[2]=14 → 36>14
  \item Step 4: i=4 → arr[4]=32 → 36>32
  \item Step 5: i=8 → overshoot → binary search 4..7 → found at index 6
\end{itemize}
\textbf{Time complexity:} O(log n), space = O(1)\\
\textbf{Advantages:} Efficient for unbounded or infinite arrays.\\
\textbf{Disadvantages:} Requires sorted array; extra step of exponential search.\\
\textbf{When to use:} Large or infinite sorted datasets.
\newpage
\section{Sorting Algorithms}

\subsection{Selection Sort}
\textbf{Idea:} Repeatedly select the smallest element and swap to front.  \\
\textbf{Steps:}
\begin{itemize}
\item Step 1: find min=5 → swap with arr[0]=34 → [5,7,23,32,34,62,36,14]

\item Step 2: min=7 → swap with arr[1]=7 → [5,7,23,32,34,62,36,14]

\item Step 3: min=14 → swap with arr[2]=23 → [5,7,14,32,34,62,36,23]

\item Step 4: min=23 → swap with arr[3]=32 → [5,7,14,23,34,62,36,32]

\item Step 5: min=32 → swap with arr[4]=34 → [5,7,14,23,32,62,36,34]

\item Step 6: min=34 → swap with arr[5]=62 → [5,7,14,23,32,34,36,62]

\item Step 7: min=36 → swap with arr[6]=36 → [5,7,14,23,32,34,36,62]
\end{itemize}
\textbf{Time complexity:} O(n²), space=O(1)\\
\textbf{Advantages:} Simple; minimal moves; good if swaps are costly.\\
\textbf{Disadvantages:} Slow for large arrays; not stable.\\
\textbf{When to use:} Small arrays, low swap cost.
\subsection{Insertion Sort}
\textbf{Idea:} Build the sorted array one element at a time. Insert each element into its correct position in the sorted part.  \\
\textbf{Steps:}  
\begin{itemize}
\item Step 1: key=7 → move 34 → insert 7 → [7,34,...]  
\item Step 2: key=23 → move 34 → insert 23 → [7,23,34,...]  
\item Step 3: key=32 → move 34 → insert 32 → [7,23,32,34,...]  
\item Step 4: key=5 → move 7,23,32,34 → insert 5 → [5,7,23,32,34,...]  
\item Step 5: key=62 → insert at end → [5,7,23,32,34,62,...]  
\item Step 6: key=36 → move 62,34 → insert 36 → [5,7,23,32,34,36,62,...]  
\item Step 7: key=14 → move 23,32,34,36,62 → insert 14 → [5,7,14,23,32,34,36,62]
\end{itemize}
\textbf{Time complexity:} O(n²), space = O(1)\\
\textbf{Advantages:} Efficient for small datasets, adaptive (fast for nearly sorted arrays)\\
\textbf{Disadvantages:} Slow for large arrays\\
\textbf{When to use:} Small arrays or mostly sorted data
\subsection{Bubble Sort}
\textbf{Idea:} Repeatedly compare adjacent elements and swap if out of order. Large values “bubble” to the end.  \\
\textbf{Steps:} 
\begin{itemize}
\item Step 1: [34, 7, 23, 32, 5, 62, 36, 14] → swap 34 & 7 → [7, 34, 23, 32, 5, 62, 36, 14] ... continue

\item Step 2: Repeat passes until no swaps → [5, 7, 14, 23, 32, 34, 36, 62]
\end{itemize}
\textbf{Time complexity:} O(n²), space = O(1)\\
\textbf{Advantages:} Simple, stable\\
\textbf{Disadvantages:} Very slow for large arrays\\
\textbf{When to use:} Only for small datasets, educational purposes
\subsection{Quick Sort}
\textbf{Idea:} Choose a pivot, partition array into smaller/larger, sort recursively  \\
\textbf{Steps (pivot = middle element 32):} 
\begin{itemize}
\item Step 1: [34, 7, 23, 32, 5, 62, 36, 14] → partition → [7, 23, 5, 14, 32, 62, 36, 34]

\item Step 2: Left [7, 23, 5, 14] pivot 23 → [7, 5, 14, 23]

\item Step 3: Right [32, 62, 36, 34] pivot 36 → [32, 34, 36, 62]

\item Step 4: Combine → [5, 7, 14, 23, 32, 34, 36, 62]
\end{itemize}
\textbf{Time complexity:} O(n log n) avg, O(n²) worst, space = O(log n)\\
\textbf{Advantages:} Very fast on average, in-place\\
\textbf{Disadvantages:} Worst case O(n²), not stable\\
\textbf{When to use:} Large arrays, general-purpose
\subsection{Heap Sort}
\textbf{Idea:} Build max heap, swap root with last element, heapify.  \\
\textbf{Steps (partial):}  
\begin{itemize}
\item Step 1: Build heap → [62, 34, 36, 32, 5, 23, 7, 14]

\item Step 2: Swap 62 with last → [14, 34, 36, 32, 5, 23, 7, 62] → heapify → [36, 32, 14, 7, 5, 23, 34, 62]

\item Step 3: Repeat → [5, 7, 14, 23, 32, 34, 36, 62]
\end{itemize}
\textbf{Time complexity:} O(n log n), space = O(1)\\
\textbf{Advantages:} In-place, guaranteed O(n log n)\\
\textbf{Disadvantages:} Not stable, more complex than simpler sorts\\
\textbf{When to use:} Large arrays when stability is not needed
\subsection{Merge Sort}
\textbf{Idea:} Divide array into halves, sort each, merge back.\\ 
\textbf{Steps:}  
\begin{itemize}
\item Step 1: Divide [34,7,23,32] and [5,62,36,14]

\item Step 2: Divide further → [34,7] [23,32] ... [5,62] [36,14]

\item Step 3: Merge → [7,34,23,32] → [7,23,32,34]

\item Step 4: Merge right → [5,62,36,14] → [5,14,36,62]

\item Step 5: Merge both → [5,7,14,23,32,34,36,62]
\end{itemize}
\textbf{Time complexity:} O(n log n), space = O(n)\\
\textbf{Advantages:} Stable, predictable O(n log n)\\
\textbf{Disadvantages:} Requires extra memory\\
\textbf{When to use:} Large arrays, stable sorting needed
\subsection{Radix Sort}
\textbf{Idea:} Sort digits from least significant to most using counting sort.  \\
\textbf{Steps (sort [34, 7, 23, 32, 5, 62, 36, 14] by units → tens):}
\begin{itemize}
\item Step 1: Units digit → [10, 5, 14, 32, 23, 36, 34, 62]

\item Step 2: Tens digit → [5, 7, 14, 23, 32, 34, 36, 62]
\end{itemize}
\textbf{Time complexity:} O(n·k), space = O(n+k)\\
\textbf{Advantages:} Very fast for integers\\
\textbf{Disadvantages:} Works only for integers or fixed-length keys\\
\textbf{When to use:} Large numbers or integers, stable
\subsection{Counting Sort}
\textbf{Idea:} Count frequency of each number → reconstruct sorted array.  \\
\textbf{Steps:} 
\begin{itemize}
\item Step 1: Count occurrences → 5:1, 7:1, 14:1, 23:1, 32:1, 34:1, 36:1, 62:1

\item Step 2: Rebuild → [5,7,14,23,32,34,36,62
\end{itemize}
\textbf{Time complexity:} O(n+k), space = O(k)\\
\textbf{Advantages:} Linear time for small range\\
\textbf{Disadvantages:} Uses extra memory for counts, only integers\\
\textbf{When to use:} Small range of integers
\subsection{Bucket Sort}
\textbf{Idea:} Divide numbers into buckets → sort each → merge.  \\
\textbf{Steps:}  
\begin{itemize}
\item Step 1: Distribute [34,7,23,32,5,62,36,14] into buckets (0-9,10-19,...)

\item Step 2: Sort each bucket → [5,7,14,23,32,34,36,62]
\end{itemize}
\textbf{Time complexity:} O(n+k), space = O(n+k)\\
\textbf{Advantages:} Efficient for uniform distribution\\
\textbf{Disadvantages:} Performance depends on distribution, extra memory\\
\textbf{When to use:} Uniformly distributed data
\subsection{Tim Sort}
\textbf{Idea:} Hybrid of merge + insertion sort (used in Python/Java).  \\
\textbf{Steps:} 
\begin{itemize}
\item Step 1: Divide array into small runs (<=32 elements) → sort with insertion sort → [5,7,14,23,32,34,36,62]

\item Step 2: Merge runs using merge sort → [5,7,14,23,32,34,36,62]
\end{itemize}
\textbf{Time complexity:} O(n log n), space = O(n)\\
\textbf{Advantages:} Stable, fast in practice, adaptive\\
\textbf{Disadvantages:} Slightly more complex implementation\\
\textbf{When to use:} General-purpose, standard library sort
\newpage
\section{Summary Table}

\begin{center}
\begin{tabular}{lccc}
\toprule
\textbf{Algorithm} & \textbf{Type} & \textbf{Time Complexity} & \textbf{Main Concept} \\
\midrule
Linear Search & Search & $O(n)$ & Check each element \\
Binary Search & Search & $O(\log n)$ & Divide by half \\
Jump Search & Search & $O(\sqrt{n})$ & Jump in blocks \\
Ternary Search & Search & $O(\log_3 n)$ & Divide into three \\
Interpolation Search & Search & $O(\log \log n)$ & Predict position \\
Exponential Search & Search & $O(\log n)$ & Expand by 2x \\
\midrule
Selection Sort & Sort & $O(n^2)$ & Pick smallest repeatedly \\
Insertion Sort & Sort & $O(n^2)$ & Insert into sorted part \\
Bubble Sort & Sort & $O(n^2)$ & Swap adjacent elements \\
Quick Sort & Sort & $O(n \log n)$ & Divide by pivot \\
Merge Sort & Sort & $O(n \log n)$ & Divide and merge \\
Heap Sort & Sort & $O(n \log n)$ & Use max heap \\
Counting Sort & Sort & $O(n + k)$ & Count frequency \\
Radix Sort & Sort & $O(n \cdot k)$ & Sort by digits \\
Bucket Sort & Sort & $O(n + k)$ & Group and merge \\
Tim Sort & Sort & $O(n \log n)$ & Hybrid (Merge + Insertion) \\
\bottomrule
\end{tabular}
\end{center}

\end{document}
